% EM structure and teaching revision 20260928
\subsection{磁镜条件与绝热近似的两重检查}
无电场且磁场随位置缓慢变化时，总动能保持不变，而第一绝热不变量
$\mu=m v_\perp^2/(2B)$ 近似保持。由初始磁场 $B_0$ 与俯仰角 $\alpha_0$ 得
\[
 v_\parallel^2=v^2\left(1-\frac{B}{B_0}\sin^2\alpha_0\right),\qquad
 B_{\mathrm{turn}}=\frac{B_0}{\sin^2\alpha_0}.
\]
磁镜反射发生在纵向速度降到零时；总速率没有下降，只是平行与垂直分量重新分配。若沿途最大场为 $B_{\max}$，反射要求 $\sin^2\alpha_0\ge B_0/B_{\max}$。
这个判据还要满足绝热条件：一回旋半径内场的相对变化很小，典型地 $r_c|\nabla B|/B\ll1$；还应检查粒子沿场前进时一回旋周期内所经历的场变化，例如 $|v_\parallel\partial_\parallel B|/(|\Omega|B)\ll1$，其中 $|\Omega|=|q|B/m$。仅有横向回旋半径很小，对 $v_\parallel\gg v_\perp$ 的粒子并不充分。若场本身随时间变化，也需回旋周期远小于场的变化时间。磁场在一个锐利界面突然改变时，应按分区轨迹匹配，不能直接宣称磁矩跨界不变。
若存在碰撞、辐射或强阻尼，总动能和磁矩可能同时改变。解题时把精确守恒量、近似不变量和纯几何约束分别列出，才知道哪一步需要额外假设。

